For \(G\sim G(n,1/2)\), with probability tending to one,
\[
 \chi_L(G)=O\left(\frac{n\log\log n}{\log n}\right)=o(n).
\]
We prove this using a random allocation of list colours
and a reserve palette. The external graph-theoretic input
is the classical ordinary-colouring bound
\(\chi(G(n,1/2))=O(n/\log n)\) with probability tending to one.
This bound, and its use in the known solution, are recorded
in Section 4 of Alon's primary paper
\emph{Choice numbers of graphs: a probabilistic approach},
\url{https://www.cs.tau.ac.il/~nogaa/PDFS/choice.pdf}.
See \url{https://www.erdosproblems.com/799}.
